\chapter{Mid-Term Examination (MTE) Model Paper}
\section*{Course: PHY175 $\cdot$ Max Marks: 50 $\cdot$ Time: 120 Minutes}

\subsection*{Part A: Short Answer Questions (10 Marks, 2 Marks Each)}
\begin{enumerate}[leftmargin=*]
    \item Define Fermi energy at absolute zero.
    \item State the difference between direct and indirect band gap semiconductors.
    \item Explain the physical significance of the Hall coefficient sign.
    \item State De Morgan's laws of Boolean algebra.
    \item Write the 2's complement representation of decimal $-25$ in 8 bits.
\end{enumerate}

\subsection*{Part B: Medium Analytical Questions (20 Marks, 5 Marks Each)}

\begin{examqbox}{1 [5 Marks]}
A 3D free electron gas at $0\,\text{K}$ has electron density $n_1$. If the density is increased to $8n_1$, prove that the Fermi energy quadruples ($E_{F2} = 4E_{F1}$).
\end{examqbox}

\begin{solbox}
Density of states integration yields $E_F = \frac{\hbar^2}{2m}(3\pi^2 n)^{2/3}$. Therefore, $\frac{E_{F2}}{E_{F1}} = \left(\frac{n_2}{n_1}\right)^{2/3} = (8)^{2/3} = (2^3)^{2/3} = 4$.
\end{solbox}

\begin{examqbox}{2 [5 Marks]}
In an intrinsic semiconductor, electron mobility is 3 times hole mobility ($\mu_e = 3\mu_h$). Calculate the percentage of total electrical conductivity contributed by electrons.
\end{examqbox}

\begin{solbox}
Total conductivity $\sigma = n_i e(\mu_e + \mu_h)$. Fractional electron conductivity is $\frac{\sigma_e}{\sigma} = \frac{n_i e \mu_e}{n_i e(\mu_e + \mu_h)} = \frac{3\mu_h}{3\mu_h + \mu_h} = \frac{3}{4} = 75\%$.
\end{solbox}

\begin{examqbox}{3 [5 Marks]}
A full-wave bridge rectifier is supplied with a $15\,\text{V}$ peak, $50\,\text{Hz}$ sinusoidal input. Diodes have $V_D = 0.7\,\text{V}$. A $1000\,\mu\text{F}$ capacitor filter and $1\,\text{k}\Omega$ load are connected. Compute peak DC voltage, ripple frequency, and peak-to-peak ripple voltage.
\end{examqbox}

\begin{solbox}
1. $V_{peak} = 15 - 2(0.7) = 13.6\,\text{V}$.\\
2. $f_{ripple} = 2 \times 50 = 100\,\text{Hz}$.\\
3. $V_{r(pp)} = \frac{V_{peak}}{2 f C R_L} = \frac{13.6}{2(50)(1000 \times 10^{-6})(1000)} = \frac{13.6}{100} = 0.136\,\text{V}$.
\end{solbox}

\begin{examqbox}{4 [5 Marks]}
Convert hexadecimal $(3\text{A}7.\text{C})_{16}$ into binary, octal, and decimal representations.
\end{examqbox}

\begin{solbox}
1. Binary: $0011\,1010\,0111.1100_2$.\\
2. Octal: $(1647.6)_8$.\\
3. Decimal: $3(256) + 10(16) + 7 + 12/16 = 768 + 160 + 7 + 0.75 = 935.75_{10}$.
\end{solbox}

\subsection*{Part C: Comprehensive Long Questions (20 Marks, 10 Marks Each)}

\begin{examqbox}{5 [10 Marks]}
(a) Derive the expression for the Hall voltage $V_H = \frac{IB}{nqt}$ and Hall coefficient $R_H$.\\
(b) A rectangular copper strip of thickness $t = 0.1\,\text{mm}$ carries $I = 10\,\text{A}$ in $B = 1.2\,\text{T}$. If $n = 8.5 \times 10^{28}\,\text{m}^{-3}$, find $V_H$.
\end{examqbox}

\begin{solbox}
\textbf{(a) Derivation:} Balancing magnetic Lorentz force $F_B = q v_d B$ with electrostatic Hall force $F_E = q E_H$ gives $E_H = v_d B$. With $v_d = \frac{I}{n q w t}$, $V_H = E_H w = \frac{I B}{n q t}$.\\
\textbf{(b) Calculation:} $V_H = \frac{(10)(1.2)}{(8.5 \times 10^{28})(1.6 \times 10^{-19})(0.1 \times 10^{-3})} = \frac{12}{1.36 \times 10^6} = 8.82\,\mu\text{V}$.
\end{solbox}

\begin{examqbox}{6 [10 Marks]}
(a) Minimize $F(A,B,C,D) = \sum m(0, 2, 5, 7, 8, 10, 13, 15)$ using a K-map.\\
(b) Explain static-1 hazards and write the hazard-free SOP expression.
\end{examqbox}

\begin{solbox}
\textbf{(a) K-Map Reduction:} Corner group $\{0, 2, 8, 10\} \implies \bar{B}\bar{D}$; Central group $\{5, 7, 13, 15\} \implies BD$. Minimized: $F = \bar{B}\bar{D} + BD$.\\
\textbf{(b) Hazard Analysis:} Since the groups are non-overlapping, input switching transitions between $B=0$ and $B=1$ can glitch. Adding consensus term $\bar{D}D=0$ is not needed here as variables are distinct, but for adjacent SOP terms, consensus loops eliminate hazards.
\end{solbox}
